\section{十九世纪战争、环境与城市网络}
内战、外战、河流改道、城市起事、军工、赈济与航运构成十九世纪的同一物质世界。本节以省际运输、城市空间、河防、工坊、港口与迁徙为线索，拒绝让战役年表遮蔽环境冲击和地方财政的长尾。
\subsection{1850—1859：区域网络、制度媒介与环境条件}
这些事件的先后可以由账本定位，但每一项影响都须回到具体区域、材料与行动者，不能被压缩为抽象的国家趋势。
\hypertarget{civic:EQ-1851-JINTIAN-UPRISING}{}\paragraph{咸丰元年（1851年）}洪秀全集团在金田起事，太平天国战争开始。本事件的战争或外交面向应与城市空间、河流道路、军需供应、避难迁徙和省级财政同时阅读，避免将军事结果抽离其物质条件。 核心取证为《清文宗实录》咸丰元年相关条；军机处档、战事方略及地方奏报。编年条目可以为行动建立相对次序，却通常不保存劳动过程、物价变动、建筑材料或非精英的叙述。账本提示的研究线索为战争动员、军需、战场暴力与地方社会研究。事件所处条件为宗教组织、地方社会冲突和清末财政治理危机交织。后续变化应限于参与者动机与人数不能由官府或太平军单方叙事概括。来源限制：以咸丰元年（1851年）的同期文书为定位起点；官修记录、奏报或外交文本服务特定机构立场，具体日期、规模、地方执行与对手视角须以独立材料复核。
\hypertarget{civic:EQ-1851-YONGAN-KINGSHIPS}{}\paragraph{咸丰元年（1851年）}太平军在永安建立王号和军事组织。本事件的战争或外交面向应与城市空间、河流道路、军需供应、避难迁徙和省级财政同时阅读，避免将军事结果抽离其物质条件。 核心取证为《清文宗实录》咸丰元年相关条；宫中朱批奏折、内阁大库题本与《清史稿》相关志传。行政文书显示的是机构如何把道路、港口、人口、宗教或案件转换为可处理的类别；没有后续县档或物质证据，不能由分类推定实际生活。账本提示的研究线索为宫廷政治、制度运作与中央—地方信息链研究。事件所处条件为起事扩展需要更稳定的指挥与合法性语言。后续变化应限于太平天国制度在流动战时环境中不断变化。来源限制：以咸丰元年（1851年）的同期文书为定位起点；官修记录、奏报或外交文本服务特定机构立场，具体日期、规模、地方执行与对手视角须以独立材料复核。
\hypertarget{civic:EQ-1852-TAIPING-NORTHWARD-MARCH}{}\paragraph{咸丰二年（1852年）}太平军北上经湖南，清军防御与地方团练开始加速组织。本事件的战争或外交面向应与城市空间、河流道路、军需供应、避难迁徙和省级财政同时阅读，避免将军事结果抽离其物质条件。 核心取证为《清文宗实录》咸丰二年相关条；军机处档、战事方略及地方奏报。编年条目可以为行动建立相对次序，却通常不保存劳动过程、物价变动、建筑材料或非精英的叙述。账本提示的研究线索为战争动员、军需、战场暴力与地方社会研究。事件所处条件为广西战场难以封堵太平军的机动路线。后续变化应限于地方武装的出现改变后续战争权力结构。来源限制：以咸丰二年（1852年）的同期文书为定位起点；官修记录、奏报或外交文本服务特定机构立场，具体日期、规模、地方执行与对手视角须以独立材料复核。
\hypertarget{civic:EQ-1852-HUNAN-ARMY-ORIGINS}{}\paragraph{咸丰二年（1852年）}曾国藩等在湖南组织团练和湘军雏形。本事件的战争或外交面向应与城市空间、河流道路、军需供应、避难迁徙和省级财政同时阅读，避免将军事结果抽离其物质条件。 核心取证为《清文宗实录》咸丰二年相关条；军机处档、战事方略及地方奏报。编年条目可以为行动建立相对次序，却通常不保存劳动过程、物价变动、建筑材料或非精英的叙述。账本提示的研究线索为战争动员、军需、战场暴力与地方社会研究。事件所处条件为八旗、绿营难以独力应对内战。后续变化应限于地方募兵和筹饷为清廷提供新力量，也提高地方军政自主性。来源限制：以咸丰二年（1852年）的同期文书为定位起点；官修记录、奏报或外交文本服务特定机构立场，具体日期、规模、地方执行与对手视角须以独立材料复核。
\hypertarget{civic:EQ-1853-NANJING-TAKING}{}\paragraph{咸丰三年（1853年）}太平军攻占南京并改名天京，建立长期政权中心。本事件的战争或外交面向应与城市空间、河流道路、军需供应、避难迁徙和省级财政同时阅读，避免将军事结果抽离其物质条件。 核心取证为《清文宗实录》咸丰三年相关条；军机处档、战事方略及地方奏报。编年条目可以为行动建立相对次序，却通常不保存劳动过程、物价变动、建筑材料或非精英的叙述。账本提示的研究线索为战争动员、军需、战场暴力与地方社会研究。事件所处条件为长江交通和清军防线崩解提供战略机会。后续变化应限于江南财政文化中心易手，战争波及全国资源网络。来源限制：以咸丰三年（1853年）的同期文书为定位起点；官修记录、奏报或外交文本服务特定机构立场，具体日期、规模、地方执行与对手视角须以独立材料复核。
\hypertarget{civic:EQ-1853-NORTHERN-EXPEDITION}{}\paragraph{咸丰三年（1853年）}太平军发动北伐，进逼华北但补给与指挥困难加重。本事件的战争或外交面向应与城市空间、河流道路、军需供应、避难迁徙和省级财政同时阅读，避免将军事结果抽离其物质条件。 核心取证为《清文宗实录》咸丰三年相关条；军机处档、战事方略及地方奏报。编年条目可以为行动建立相对次序，却通常不保存劳动过程、物价变动、建筑材料或非精英的叙述。账本提示的研究线索为战争动员、军需、战场暴力与地方社会研究。事件所处条件为天京建立后希望扩大战略压力。后续变化应限于北伐失败限制太平天国对北方的直接威胁。来源限制：以咸丰三年（1853年）的同期文书为定位起点；官修记录、奏报或外交文本服务特定机构立场，具体日期、规模、地方执行与对手视角须以独立材料复核。
\hypertarget{civic:EQ-1853-WESTERN-EXPEDITION}{}\paragraph{咸丰三年（1853年）}太平军西征争夺长江中上游，战事牵动安徽、江西、湖北。本事件的战争或外交面向应与城市空间、河流道路、军需供应、避难迁徙和省级财政同时阅读，避免将军事结果抽离其物质条件。 核心取证为《清文宗实录》咸丰三年相关条；军机处档、战事方略及地方奏报。编年条目可以为行动建立相对次序，却通常不保存劳动过程、物价变动、建筑材料或非精英的叙述。账本提示的研究线索为战争动员、军需、战场暴力与地方社会研究。事件所处条件为天京需要控制粮食、兵员和长江交通。后续变化应限于区域战场长期化，地方社会承受征发与迁徙。来源限制：以咸丰三年（1853年）的同期文书为定位起点；官修记录、奏报或外交文本服务特定机构立场，具体日期、规模、地方执行与对手视角须以独立材料复核。
\hypertarget{civic:EQ-1853-SMALL-SWORDS-SHANGHAI}{}\paragraph{咸丰三年（1853年）}上海小刀会起事，租界和华人城市空间关系发生变化。本事件的战争或外交面向应与城市空间、河流道路、军需供应、避难迁徙和省级财政同时阅读，避免将军事结果抽离其物质条件。 核心取证为《清文宗实录》咸丰三年相关条；地方志、督抚奏折及地方档案。编年条目可以为行动建立相对次序，却通常不保存劳动过程、物价变动、建筑材料或非精英的叙述。账本提示的研究线索为地方档案、迁徙、宗教、族群与社会动员研究。事件所处条件为城市商业、秘密结社和地方治理矛盾交织。后续变化应限于租界保护与地方镇压的边界需结合中外档案分析。来源限制：以咸丰三年（1853年）的同期文书为定位起点；官修记录、奏报或外交文本服务特定机构立场，具体日期、规模、地方执行与对手视角须以独立材料复核。
\hypertarget{civic:EQ-1854-TAIPING-NORTHERN-DEFEAT}{}\paragraph{咸丰四年（1854年）}太平军北伐受挫，清军逐步恢复华北防线。本事件的战争或外交面向应与城市空间、河流道路、军需供应、避难迁徙和省级财政同时阅读，避免将军事结果抽离其物质条件。 核心取证为《清文宗实录》咸丰四年相关条；军机处档、战事方略及地方奏报。编年条目可以为行动建立相对次序，却通常不保存劳动过程、物价变动、建筑材料或非精英的叙述。账本提示的研究线索为战争动员、军需、战场暴力与地方社会研究。事件所处条件为长距离行军、冬季补给与地方支持不足限制北伐。后续变化应限于太平天国转而依赖长江流域与地方战场。来源限制：以咸丰四年（1854年）的同期文书为定位起点；官修记录、奏报或外交文本服务特定机构立场，具体日期、规模、地方执行与对手视角须以独立材料复核。
\hypertarget{civic:EQ-1854-XIANG-ARMY-ORGANIZATION}{}\paragraph{咸丰四年（1854年）}湘军编制、饷源和将领网络逐步成形。本事件的战争或外交面向应与城市空间、河流道路、军需供应、避难迁徙和省级财政同时阅读，避免将军事结果抽离其物质条件。 核心取证为《清文宗实录》咸丰四年相关条；宫中朱批奏折、内阁大库题本与《清史稿》相关志传。行政文书显示的是机构如何把道路、港口、人口、宗教或案件转换为可处理的类别；没有后续县档或物质证据，不能由分类推定实际生活。账本提示的研究线索为宫廷政治、制度运作与中央—地方信息链研究。事件所处条件为清廷需要可用的地方军事力量。后续变化应限于军事私人网络与国家授权之间的界线长期模糊。来源限制：以咸丰四年（1854年）的同期文书为定位起点；官修记录、奏报或外交文本服务特定机构立场，具体日期、规模、地方执行与对手视角须以独立材料复核。
\hypertarget{civic:EQ-1855-YELLOW-RIVER-COURSE-SHIFT}{}\paragraph{咸丰五年（1855年）}黄河在铜瓦厢决口改道北流，重塑华北水系与社会经济。本事件的战争或外交面向应与城市空间、河流道路、军需供应、避难迁徙和省级财政同时阅读，避免将军事结果抽离其物质条件。 核心取证为《清文宗实录》咸丰五年相关条；地方志、督抚奏折及地方档案。编年条目可以为行动建立相对次序，却通常不保存劳动过程、物价变动、建筑材料或非精英的叙述。账本提示的研究线索为地方档案、迁徙、宗教、族群与社会动员研究。事件所处条件为长期河工压力和水文变化共同作用。后续变化应限于河道变迁的时间、影响范围和人口后果须依地方资料分区考察。来源限制：以咸丰五年（1855年）的同期文书为定位起点；官修记录、奏报或外交文本服务特定机构立场，具体日期、规模、地方执行与对手视角须以独立材料复核。
\hypertarget{civic:EQ-1855-NIAN-REBELLION-EXPANSION}{}\paragraph{咸丰五年（1855年）}捻军在华北、安徽等地活动扩大，形成与太平战争交织的战场。本事件的战争或外交面向应与城市空间、河流道路、军需供应、避难迁徙和省级财政同时阅读，避免将军事结果抽离其物质条件。 核心取证为《清文宗实录》咸丰五年相关条；军机处档、战事方略及地方奏报。编年条目可以为行动建立相对次序，却通常不保存劳动过程、物价变动、建筑材料或非精英的叙述。账本提示的研究线索为战争动员、军需、战场暴力与地方社会研究。事件所处条件为灾荒、流民、盐枭和地方武装网络提供条件。后续变化应限于捻军组织结构和社会基础不应由“匪”一词简化。来源限制：以咸丰五年（1855年）的同期文书为定位起点；官修记录、奏报或外交文本服务特定机构立场，具体日期、规模、地方执行与对手视角须以独立材料复核。
\hypertarget{civic:EQ-1856-TIANJING-INCIDENT}{}\paragraph{咸丰六年（1856年）}天京事变中太平天国领导层发生内讧和大规模杀戮。本事件的战争或外交面向应与城市空间、河流道路、军需供应、避难迁徙和省级财政同时阅读，避免将军事结果抽离其物质条件。 核心取证为《清文宗实录》咸丰六年相关条；宫中朱批奏折、内阁大库题本与《清史稿》相关志传。行政文书显示的是机构如何把道路、港口、人口、宗教或案件转换为可处理的类别；没有后续县档或物质证据，不能由分类推定实际生活。账本提示的研究线索为宫廷政治、制度运作与中央—地方信息链研究。事件所处条件为权力集中、军政分工与宗教权威冲突累积。后续变化应限于事件严重削弱太平天国协调能力，细节和死亡数存在争议。来源限制：以咸丰六年（1856年）的同期文书为定位起点；官修记录、奏报或外交文本服务特定机构立场，具体日期、规模、地方执行与对手视角须以独立材料复核。
\hypertarget{civic:EQ-1856-SECOND-OPIUM-WAR-BEGIN}{}\paragraph{咸丰六年（1856年）}亚罗号事件等冲突后英法对清开战，第二次鸦片战争开始。本事件的战争或外交面向应与城市空间、河流道路、军需供应、避难迁徙和省级财政同时阅读，避免将军事结果抽离其物质条件。 核心取证为《清文宗实录》咸丰六年相关条；军机处档、战事方略及地方奏报。编年条目可以为行动建立相对次序，却通常不保存劳动过程、物价变动、建筑材料或非精英的叙述。账本提示的研究线索为战争动员、军需、战场暴力与地方社会研究。事件所处条件为条约修订、领事权与商业扩张要求升级。后续变化应限于战争与内战并行，迫使清廷分散军事和财政资源。来源限制：以咸丰六年（1856年）的同期文书为定位起点；官修记录、奏报或外交文本服务特定机构立场，具体日期、规模、地方执行与对手视角须以独立材料复核。
\hypertarget{civic:EQ-1857-GUANGZHOU-FALL}{}\paragraph{咸丰七年（1857年）}英法联军攻占广州，叶名琛被俘。本事件的战争或外交面向应与城市空间、河流道路、军需供应、避难迁徙和省级财政同时阅读，避免将军事结果抽离其物质条件。 核心取证为《清文宗实录》咸丰七年相关条；军机处档、战事方略及地方奏报。编年条目可以为行动建立相对次序，却通常不保存劳动过程、物价变动、建筑材料或非精英的叙述。账本提示的研究线索为战争动员、军需、战场暴力与地方社会研究。事件所处条件为珠江口海军优势和清军防御薄弱。后续变化应限于广东地方治理、商贸和对外关系受到持续冲击。来源限制：以咸丰七年（1857年）的同期文书为定位起点；官修记录、奏报或外交文本服务特定机构立场，具体日期、规模、地方执行与对手视角须以独立材料复核。
\hypertarget{civic:EQ-1858-TIANJIN-TREATIES}{}\paragraph{咸丰八年（1858年）}《天津条约》等在战争压力下签订，扩大通商、使节与传教权利。本事件的战争或外交面向应与城市空间、河流道路、军需供应、避难迁徙和省级财政同时阅读，避免将军事结果抽离其物质条件。 核心取证为《清文宗实录》咸丰八年相关条；《筹办夷务始末》相关年卷与中外照会、条约文本。行政文书显示的是机构如何把道路、港口、人口、宗教或案件转换为可处理的类别；没有后续县档或物质证据，不能由分类推定实际生活。账本提示的研究线索为条约文本、外交档案、口岸社会与国际法实践研究。事件所处条件为联军北上威胁京师迫使清廷议约。后续变化应限于签署不等于立即批准或实施，后续冲突仍在发生。来源限制：以咸丰八年（1858年）的同期文书为定位起点；官修记录、奏报或外交文本服务特定机构立场，具体日期、规模、地方执行与对手视角须以独立材料复核。
\hypertarget{civic:EQ-1859-TAKU-BATTLE}{}\paragraph{咸丰九年（1859年）}大沽口战事使英法换约受阻，战争再度升级。本事件的战争或外交面向应与城市空间、河流道路、军需供应、避难迁徙和省级财政同时阅读，避免将军事结果抽离其物质条件。 核心取证为《清文宗实录》咸丰九年相关条；军机处档、战事方略及地方奏报。编年条目可以为行动建立相对次序，却通常不保存劳动过程、物价变动、建筑材料或非精英的叙述。账本提示的研究线索为战争动员、军需、战场暴力与地方社会研究。事件所处条件为条约批准和使节入京问题未获解决。后续变化应限于清方短期胜利未改变总体军力与外交劣势。来源限制：以咸丰九年（1859年）的同期文书为定位起点；官修记录、奏报或外交文本服务特定机构立场，具体日期、规模、地方执行与对手视角须以独立材料复核。
\subsection{1860—1869：区域网络、制度媒介与环境条件}
这些事件的先后可以由账本定位，但每一项影响都须回到具体区域、材料与行动者，不能被压缩为抽象的国家趋势。
\hypertarget{civic:EQ-1860-BEIJING-CAMPAIGN}{}\paragraph{咸丰十年（1860年）}英法联军进逼北京，圆明园遭焚掠，咸丰帝出走热河。本事件的战争或外交面向应与城市空间、河流道路、军需供应、避难迁徙和省级财政同时阅读，避免将军事结果抽离其物质条件。 核心取证为《清文宗实录》咸丰十年相关条；军机处档、战事方略及地方奏报。编年条目可以为行动建立相对次序，却通常不保存劳动过程、物价变动、建筑材料或非精英的叙述。账本提示的研究线索为战争动员、军需、战场暴力与地方社会研究。事件所处条件为大沽后战争升级与清军防线崩溃。后续变化应限于战争记忆、掠夺和城市影响需结合中外见证与物质证据。来源限制：以咸丰十年（1860年）的同期文书为定位起点；官修记录、奏报或外交文本服务特定机构立场，具体日期、规模、地方执行与对手视角须以独立材料复核。
\hypertarget{civic:EQ-1860-CONVENTION-OF-PEKING}{}\paragraph{咸丰十年（1860年）}《北京条约》及相关中俄条约确认和扩大列强权益并涉及东北边界。本事件的战争或外交面向应与城市空间、河流道路、军需供应、避难迁徙和省级财政同时阅读，避免将军事结果抽离其物质条件。 核心取证为《清文宗实录》咸丰十年相关条；《筹办夷务始末》相关年卷与中外照会、条约文本。行政文书显示的是机构如何把道路、港口、人口、宗教或案件转换为可处理的类别；没有后续县档或物质证据，不能由分类推定实际生活。账本提示的研究线索为条约文本、外交档案、口岸社会与国际法实践研究。事件所处条件为北京陷落迫使清廷接受新安排。后续变化应限于各条约文本、边界实施和土地控制不可混为一谈。来源限制：以咸丰十年（1860年）的同期文书为定位起点；官修记录、奏报或外交文本服务特定机构立场，具体日期、规模、地方执行与对手视角须以独立材料复核。
\hypertarget{civic:EQ-1861-XIANFENG-DEATH}{}\paragraph{咸丰十一年（1861年）}咸丰帝去世，同治帝即位，顾命大臣与两宫太后权力冲突爆发。本事件的战争或外交面向应与城市空间、河流道路、军需供应、避难迁徙和省级财政同时阅读，避免将军事结果抽离其物质条件。 核心取证为《清文宗实录》咸丰十一年相关条；宫中朱批奏折、内阁大库题本与《清史稿》相关志传。行政文书显示的是机构如何把道路、港口、人口、宗教或案件转换为可处理的类别；没有后续县档或物质证据，不能由分类推定实际生活。账本提示的研究线索为宫廷政治、制度运作与中央—地方信息链研究。事件所处条件为战时政局与幼主继承形成权力真空。后续变化应限于辛酉政变重塑清廷中枢决策结构。来源限制：以咸丰十一年（1861年）的同期文书为定位起点；官修记录、奏报或外交文本服务特定机构立场，具体日期、规模、地方执行与对手视角须以独立材料复核。
\hypertarget{civic:EQ-1861-XINYOU-COUP}{}\paragraph{咸丰十一年（1861年）}慈禧、慈安与恭亲王发动辛酉政变，肃顺等顾命大臣失势。本事件的战争或外交面向应与城市空间、河流道路、军需供应、避难迁徙和省级财政同时阅读，避免将军事结果抽离其物质条件。 核心取证为《清文宗实录》咸丰十一年相关条；宫中朱批奏折、内阁大库题本与《清史稿》相关志传。行政文书显示的是机构如何把道路、港口、人口、宗教或案件转换为可处理的类别；没有后续县档或物质证据，不能由分类推定实际生活。账本提示的研究线索为宫廷政治、制度运作与中央—地方信息链研究。事件所处条件为幼主摄政安排与京师政治网络冲突。后续变化应限于两宫垂帘与恭亲王参与中枢成为同治初年格局。来源限制：以咸丰十一年（1861年）的同期文书为定位起点；官修记录、奏报或外交文本服务特定机构立场，具体日期、规模、地方执行与对手视角须以独立材料复核。
\hypertarget{civic:EQ-1861-ZONGLI-YAMEN}{}\paragraph{咸丰十一年（1861年）}总理各国事务衙门设立，清廷开始形成常设近代外交处理机构。本事件的战争或外交面向应与城市空间、河流道路、军需供应、避难迁徙和省级财政同时阅读，避免将军事结果抽离其物质条件。 核心取证为《清文宗实录》咸丰十一年相关条；《筹办夷务始末》相关年卷与中外照会、条约文本。行政文书显示的是机构如何把道路、港口、人口、宗教或案件转换为可处理的类别；没有后续县档或物质证据，不能由分类推定实际生活。账本提示的研究线索为条约文本、外交档案、口岸社会与国际法实践研究。事件所处条件为条约体系和驻京使节要求持续对外交涉能力。后续变化应限于机构职责不断调整，不能视为完整现代外交部。来源限制：以咸丰十一年（1861年）的同期文书为定位起点；官修记录、奏报或外交文本服务特定机构立场，具体日期、规模、地方执行与对手视角须以独立材料复核。
\hypertarget{civic:EQ-1862-SELF-STRENGTHENING-ORIGINS}{}\paragraph{同治元年（1862年）}洋务实践在军械、翻译、外交和地方军政中展开。本事件的战争或外交面向应与城市空间、河流道路、军需供应、避难迁徙和省级财政同时阅读，避免将军事结果抽离其物质条件。 核心取证为《清穆宗实录》同治元年相关条；宫中朱批奏折、内阁大库题本与《清史稿》相关志传。行政文书显示的是机构如何把道路、港口、人口、宗教或案件转换为可处理的类别；没有后续县档或物质证据，不能由分类推定实际生活。账本提示的研究线索为宫廷政治、制度运作与中央—地方信息链研究。事件所处条件为内战与列强压力促使官员寻求技术和制度调整。后续变化应限于不同地区的洋务项目具有不同资金、目标和绩效。来源限制：以同治元年（1862年）的同期文书为定位起点；官修记录、奏报或外交文本服务特定机构立场，具体日期、规模、地方执行与对手视角须以独立材料复核。
\hypertarget{civic:EQ-1862-SHANGHAI-FOREIGN-ARMS}{}\paragraph{同治元年（1862年）}上海的军火、雇佣力量和租界网络影响清军对太平军作战。本事件的战争或外交面向应与城市空间、河流道路、军需供应、避难迁徙和省级财政同时阅读，避免将军事结果抽离其物质条件。 核心取证为《清穆宗实录》同治元年相关条；军机处档、战事方略及地方奏报。编年条目可以为行动建立相对次序，却通常不保存劳动过程、物价变动、建筑材料或非精英的叙述。账本提示的研究线索为战争动员、军需、战场暴力与地方社会研究。事件所处条件为江南战场接近条约口岸与国际商业网络。后续变化应限于外来军事技术和人员角色须避免夸大为决定性单因。来源限制：以同治元年（1862年）的同期文书为定位起点；官修记录、奏报或外交文本服务特定机构立场，具体日期、规模、地方执行与对手视角须以独立材料复核。
\hypertarget{civic:EQ-1863-EVER-VICTORIOUS-ARMY}{}\paragraph{同治二年（1863年）}常胜军在江南战场参与清军作战，戈登接掌后扩大影响。本事件的战争或外交面向应与城市空间、河流道路、军需供应、避难迁徙和省级财政同时阅读，避免将军事结果抽离其物质条件。 核心取证为《清穆宗实录》同治二年相关条；军机处档、战事方略及地方奏报。编年条目可以为行动建立相对次序，却通常不保存劳动过程、物价变动、建筑材料或非精英的叙述。账本提示的研究线索为战争动员、军需、战场暴力与地方社会研究。事件所处条件为上海周边的财政、外商和地方军政资源可被军事化。后续变化应限于常胜军规模有限，其作用需放入湘淮军和地方战场比较。来源限制：以同治二年（1863年）的同期文书为定位起点；官修记录、奏报或外交文本服务特定机构立场，具体日期、规模、地方执行与对手视角须以独立材料复核。
\hypertarget{civic:EQ-1863-NIAN-WAR-MOBILIZATION}{}\paragraph{同治二年（1863年）}清廷调集湘淮军和地方力量应对捻军，华北战事加剧。本事件的战争或外交面向应与城市空间、河流道路、军需供应、避难迁徙和省级财政同时阅读，避免将军事结果抽离其物质条件。 核心取证为《清穆宗实录》同治二年相关条；军机处档、战事方略及地方奏报。编年条目可以为行动建立相对次序，却通常不保存劳动过程、物价变动、建筑材料或非精英的叙述。账本提示的研究线索为战争动员、军需、战场暴力与地方社会研究。事件所处条件为太平战场与北方流动战场同时消耗军费。后续变化应限于军队调动改变省际财政与地方权力关系。来源限制：以同治二年（1863年）的同期文书为定位起点；官修记录、奏报或外交文本服务特定机构立场，具体日期、规模、地方执行与对手视角须以独立材料复核。
\hypertarget{civic:EQ-1864-TIANJING-FALL}{}\paragraph{同治三年（1864年）}湘军攻陷天京，太平天国核心政权覆亡。本事件的战争或外交面向应与城市空间、河流道路、军需供应、避难迁徙和省级财政同时阅读，避免将军事结果抽离其物质条件。 核心取证为《清穆宗实录》同治三年相关条；军机处档、战事方略及地方奏报。编年条目可以为行动建立相对次序，却通常不保存劳动过程、物价变动、建筑材料或非精英的叙述。账本提示的研究线索为战争动员、军需、战场暴力与地方社会研究。事件所处条件为长期围困、粮食危机和内部裂解削弱天京。后续变化应限于战后屠杀、人口损失和城市重建须以多种材料谨慎处理。来源限制：以同治三年（1864年）的同期文书为定位起点；官修记录、奏报或外交文本服务特定机构立场，具体日期、规模、地方执行与对手视角须以独立材料复核。
\hypertarget{civic:EQ-1865-DUNGAN-REVOLT}{}\paragraph{同治四年（1865年）}西北回民起事在陕西、甘肃等地扩展，形成长期战争。本事件的战争或外交面向应与城市空间、河流道路、军需供应、避难迁徙和省级财政同时阅读，避免将军事结果抽离其物质条件。 核心取证为《清穆宗实录》同治四年相关条；军机处档、战事方略及地方奏报。编年条目可以为行动建立相对次序，却通常不保存劳动过程、物价变动、建筑材料或非精英的叙述。账本提示的研究线索为战争动员、军需、战场暴力与地方社会研究。事件所处条件为地方冲突、军政失序和族群分类相互激化。后续变化应限于动机、伤亡和社区边界须用地方与多语材料审慎叙述。来源限制：以同治四年（1865年）的同期文书为定位起点；官修记录、奏报或外交文本服务特定机构立场，具体日期、规模、地方执行与对手视角须以独立材料复核。
\hypertarget{civic:EQ-1865-NIAN-LEADER-DEATH}{}\paragraph{同治四年（1865年）}张乐行等捻军领导人先后失势或被杀，捻军仍持续活动。本事件的战争或外交面向应与城市空间、河流道路、军需供应、避难迁徙和省级财政同时阅读，避免将军事结果抽离其物质条件。 核心取证为《清穆宗实录》同治四年相关条；军机处档、战事方略及地方奏报。编年条目可以为行动建立相对次序，却通常不保存劳动过程、物价变动、建筑材料或非精英的叙述。账本提示的研究线索为战争动员、军需、战场暴力与地方社会研究。事件所处条件为清军加强围堵与地方武装协作。后续变化应限于领导层变化不等于流动作战网络立即瓦解。来源限制：以同治四年（1865年）的同期文书为定位起点；官修记录、奏报或外交文本服务特定机构立场，具体日期、规模、地方执行与对手视角须以独立材料复核。
\hypertarget{civic:EQ-1866-FUJIAN-SHIPYARD}{}\paragraph{同治五年（1866年）}福州船政局创办，成为洋务造船、教育与技术翻译的重要项目。本事件的战争或外交面向应与城市空间、河流道路、军需供应、避难迁徙和省级财政同时阅读，避免将军事结果抽离其物质条件。 核心取证为《清穆宗实录》同治五年相关条；户部奏销、盐政、河工或海关档案。行政文书显示的是机构如何把道路、港口、人口、宗教或案件转换为可处理的类别；没有后续县档或物质证据，不能由分类推定实际生活。账本提示的研究线索为财政账册、市场网络、灾害环境与区域经济研究。事件所处条件为海防压力和技术学习需求增长。后续变化应限于厂局经费、产能和军事效果需用账册与具体船只记录评估。来源限制：以同治五年（1866年）的同期文书为定位起点；官修记录、奏报或外交文本服务特定机构立场，具体日期、规模、地方执行与对手视角须以独立材料复核。
\hypertarget{civic:EQ-1866-NINGBO-SHANGHAI-INDUSTRY}{}\paragraph{同治五年（1866年）}江南制造局等军工项目在战后江南环境中继续发展。本事件的战争或外交面向应与城市空间、河流道路、军需供应、避难迁徙和省级财政同时阅读，避免将军事结果抽离其物质条件。 核心取证为《清穆宗实录》同治五年相关条；户部奏销、盐政、河工或海关档案。行政文书显示的是机构如何把道路、港口、人口、宗教或案件转换为可处理的类别；没有后续县档或物质证据，不能由分类推定实际生活。账本提示的研究线索为财政账册、市场网络、灾害环境与区域经济研究。事件所处条件为湘淮军和条约口岸带来技术、资金与人才条件。后续变化应限于“自强”口号不能替代对生产规模和依赖进口的检验。来源限制：以同治五年（1866年）的同期文书为定位起点；官修记录、奏报或外交文本服务特定机构立场，具体日期、规模、地方执行与对手视角须以独立材料复核。
\hypertarget{civic:EQ-1867-NIAN-CAVALRY-CAMPAIGN}{}\paragraph{同治六年（1867年）}清军以河防、堡垒和机动军队围堵捻军骑兵。本事件的战争或外交面向应与城市空间、河流道路、军需供应、避难迁徙和省级财政同时阅读，避免将军事结果抽离其物质条件。 核心取证为《清穆宗实录》同治六年相关条；军机处档、战事方略及地方奏报。编年条目可以为行动建立相对次序，却通常不保存劳动过程、物价变动、建筑材料或非精英的叙述。账本提示的研究线索为战争动员、军需、战场暴力与地方社会研究。事件所处条件为捻军流动作战突破传统省界防御。后续变化应限于军事策略的效果应与地方破坏和民众负担并看。来源限制：以同治六年（1867年）的同期文书为定位起点；官修记录、奏报或外交文本服务特定机构立场，具体日期、规模、地方执行与对手视角须以独立材料复核。
\hypertarget{civic:EQ-1868-NIAN-SUPPRESSION}{}\paragraph{同治七年（1868年）}东捻军等主力被击败，捻军大规模战争告终。本事件的战争或外交面向应与城市空间、河流道路、军需供应、避难迁徙和省级财政同时阅读，避免将军事结果抽离其物质条件。 核心取证为《清穆宗实录》同治七年相关条；军机处档、战事方略及地方奏报。编年条目可以为行动建立相对次序，却通常不保存劳动过程、物价变动、建筑材料或非精英的叙述。账本提示的研究线索为战争动员、军需、战场暴力与地方社会研究。事件所处条件为淮军等部队的围堵和战场分割发挥作用。后续变化应限于战后地方武装和财政动员遗产继续存在。来源限制：以同治七年（1868年）的同期文书为定位起点；官修记录、奏报或外交文本服务特定机构立场，具体日期、规模、地方执行与对手视角须以独立材料复核。
\hypertarget{civic:EQ-1869-YANGZHOU-SALT-CRISIS}{}\paragraph{同治八年（1869年）}两淮盐政改革、商欠与财政重整成为中兴时期的重要议题。本事件的战争或外交面向应与城市空间、河流道路、军需供应、避难迁徙和省级财政同时阅读，避免将军事结果抽离其物质条件。 核心取证为《清穆宗实录》同治八年相关条；户部奏销、盐政、河工或海关档案。行政文书显示的是机构如何把道路、港口、人口、宗教或案件转换为可处理的类别；没有后续县档或物质证据，不能由分类推定实际生活。账本提示的研究线索为财政账册、市场网络、灾害环境与区域经济研究。事件所处条件为战乱破坏盐业网络，厘金与新税制改变利益结构。后续变化应限于改革效果需比较名义制度、实收与商人账目。来源限制：以同治八年（1869年）的同期文书为定位起点；官修记录、奏报或外交文本服务特定机构立场，具体日期、规模、地方执行与对手视角须以独立材料复核。
\subsection{1870—1879：区域网络、制度媒介与环境条件}
这些事件的先后可以由账本定位，但每一项影响都须回到具体区域、材料与行动者，不能被压缩为抽象的国家趋势。
\hypertarget{civic:EQ-1870-TIANJIN-MASSACRE}{}\paragraph{同治九年（1870年）}天津教案发生，法国领事、修女和中国民众死伤，引发外交危机。本事件的战争或外交面向应与城市空间、河流道路、军需供应、避难迁徙和省级财政同时阅读，避免将军事结果抽离其物质条件。 核心取证为《清穆宗实录》同治九年相关条；《筹办夷务始末》相关年卷与中外照会、条约文本。行政文书显示的是机构如何把道路、港口、人口、宗教或案件转换为可处理的类别；没有后续县档或物质证据，不能由分类推定实际生活。账本提示的研究线索为条约文本、外交档案、口岸社会与国际法实践研究。事件所处条件为教会活动、地方谣言和条约保护关系高度紧张。后续变化应限于责任叙述在中法档案中差异显著，须避免单一归因。来源限制：以同治九年（1870年）的同期文书为定位起点；官修记录、奏报或外交文本服务特定机构立场，具体日期、规模、地方执行与对手视角须以独立材料复核。
\hypertarget{civic:EQ-1870-ZHILI-FAMINE-RELIEF}{}\paragraph{同治九年（1870年）}华北灾荒和赈务问题促使官员、士绅和国际救济网络介入。本事件的战争或外交面向应与城市空间、河流道路、军需供应、避难迁徙和省级财政同时阅读，避免将军事结果抽离其物质条件。 核心取证为《清穆宗实录》同治九年相关条；地方志、督抚奏折及地方档案。编年条目可以为行动建立相对次序，却通常不保存劳动过程、物价变动、建筑材料或非精英的叙述。账本提示的研究线索为地方档案、迁徙、宗教、族群与社会动员研究。事件所处条件为旱灾、粮价、交通和地方储备共同影响灾情。后续变化应限于慈善报告与官府奏报的覆盖范围和数字需核对。来源限制：以同治九年（1870年）的同期文书为定位起点；官修记录、奏报或外交文本服务特定机构立场，具体日期、规模、地方执行与对手视角须以独立材料复核。
\hypertarget{civic:EQ-1871-MUDAN-INCIDENT}{}\paragraph{同治十年（1871年）}琉球船民在台湾牡丹社一带遇害，事件成为日本对台政策借口。本事件的战争或外交面向应与城市空间、河流道路、军需供应、避难迁徙和省级财政同时阅读，避免将军事结果抽离其物质条件。 核心取证为《清穆宗实录》同治十年相关条；《筹办夷务始末》相关年卷与中外照会、条约文本。行政文书显示的是机构如何把道路、港口、人口、宗教或案件转换为可处理的类别；没有后续县档或物质证据，不能由分类推定实际生活。账本提示的研究线索为条约文本、外交档案、口岸社会与国际法实践研究。事件所处条件为海难、原住民地域与清朝地方治理边界交错。后续变化应限于事件叙述须区分原住民行动、清方管辖与日方外交利用。来源限制：以同治十年（1871年）的同期文书为定位起点；官修记录、奏报或外交文本服务特定机构立场，具体日期、规模、地方执行与对手视角须以独立材料复核。
\hypertarget{civic:EQ-1871-QING-STUDENTS-ABROAD}{}\paragraph{同治十年（1871年）}清廷派遣幼童赴美留学的计划启动，洋务教育进入海外培训阶段。本事件的战争或外交面向应与城市空间、河流道路、军需供应、避难迁徙和省级财政同时阅读，避免将军事结果抽离其物质条件。 核心取证为《清穆宗实录》同治十年相关条；内阁档、文集或报刊相关记录。编年条目可以为行动建立相对次序，却通常不保存劳动过程、物价变动、建筑材料或非精英的叙述。账本提示的研究线索为知识生产、教育、宗教与公共传播研究。事件所处条件为技术、语言与外交人才需求上升。后续变化应限于选派和留学经历不代表全国教育制度已转型。来源限制：以同治十年（1871年）的同期文书为定位起点；官修记录、奏报或外交文本服务特定机构立场，具体日期、规模、地方执行与对手视角须以独立材料复核。
\hypertarget{civic:EQ-1872-CHINA-MERCHANTS-STEAM-NAVIGATION}{}\paragraph{同治十一年（1872年）}轮船招商局创办，试图以官督商办方式发展近代航运。本事件的战争或外交面向应与城市空间、河流道路、军需供应、避难迁徙和省级财政同时阅读，避免将军事结果抽离其物质条件。 核心取证为《清穆宗实录》同治十一年相关条；户部奏销、盐政、河工或海关档案。行政文书显示的是机构如何把道路、港口、人口、宗教或案件转换为可处理的类别；没有后续县档或物质证据，不能由分类推定实际生活。账本提示的研究线索为财政账册、市场网络、灾害环境与区域经济研究。事件所处条件为外轮竞争和沿海、长江运输需求推动新企业。后续变化应限于企业性质、资本来源和经营绩效需依账册分析。来源限制：以同治十一年（1872年）的同期文书为定位起点；官修记录、奏报或外交文本服务特定机构立场，具体日期、规模、地方执行与对手视角须以独立材料复核。
\hypertarget{civic:EQ-1873-TONGZHI-EMPEROR-REIGN}{}\paragraph{同治十二年（1873年）}同治帝亲政，垂帘体制与皇帝个人政治角色发生调整。本事件的战争或外交面向应与城市空间、河流道路、军需供应、避难迁徙和省级财政同时阅读，避免将军事结果抽离其物质条件。 核心取证为《清穆宗实录》同治十二年相关条；宫中朱批奏折、内阁大库题本与《清史稿》相关志传。行政文书显示的是机构如何把道路、港口、人口、宗教或案件转换为可处理的类别；没有后续县档或物质证据，不能由分类推定实际生活。账本提示的研究线索为宫廷政治、制度运作与中央—地方信息链研究。事件所处条件为幼帝成年改变摄政安排。后续变化应限于实际决策仍受太后、军机和重臣网络制约。来源限制：以同治十二年（1873年）的同期文书为定位起点；官修记录、奏报或外交文本服务特定机构立场，具体日期、规模、地方执行与对手视角须以独立材料复核。
\hypertarget{civic:EQ-1873-JAPAN-RYUKYU-DISPUTE}{}\paragraph{同治十二年（1873年）}日本围绕琉球与台湾事务加强对清交涉。本事件的战争或外交面向应与城市空间、河流道路、军需供应、避难迁徙和省级财政同时阅读，避免将军事结果抽离其物质条件。 核心取证为《清穆宗实录》同治十二年相关条；《筹办夷务始末》相关年卷与中外照会、条约文本。行政文书显示的是机构如何把道路、港口、人口、宗教或案件转换为可处理的类别；没有后续县档或物质证据，不能由分类推定实际生活。账本提示的研究线索为条约文本、外交档案、口岸社会与国际法实践研究。事件所处条件为日本明治国家扩张与东亚朝贡关系重组。后续变化应限于册封、属国和主权词汇应置于各方文本语境中。来源限制：以同治十二年（1873年）的同期文书为定位起点；官修记录、奏报或外交文本服务特定机构立场，具体日期、规模、地方执行与对手视角须以独立材料复核。
\hypertarget{civic:EQ-1874-JAPANESE-TAIWAN-EXPEDITION}{}\paragraph{同治十三年（1874年）}日本以牡丹社事件为由出兵台湾南部，清廷被迫应对。本事件的战争或外交面向应与城市空间、河流道路、军需供应、避难迁徙和省级财政同时阅读，避免将军事结果抽离其物质条件。 核心取证为《清穆宗实录》同治十三年相关条；军机处档、战事方略及地方奏报。编年条目可以为行动建立相对次序，却通常不保存劳动过程、物价变动、建筑材料或非精英的叙述。账本提示的研究线索为战争动员、军需、战场暴力与地方社会研究。事件所处条件为清廷对台湾东部和原住民地区的治理能力有限。后续变化应限于战后赔款与外交安排刺激清廷加强台湾经营。来源限制：以同治十三年（1874年）的同期文书为定位起点；官修记录、奏报或外交文本服务特定机构立场，具体日期、规模、地方执行与对手视角须以独立材料复核。
\hypertarget{civic:EQ-1874-TONGZHI-DEATH}{}\paragraph{同治十三年（1874年）}同治帝去世，慈禧选择光绪帝继位，再次形成幼主政治。本事件的战争或外交面向应与城市空间、河流道路、军需供应、避难迁徙和省级财政同时阅读，避免将军事结果抽离其物质条件。 核心取证为《清穆宗实录》同治十三年相关条；宫中朱批奏折、内阁大库题本与《清史稿》相关志传。行政文书显示的是机构如何把道路、港口、人口、宗教或案件转换为可处理的类别；没有后续县档或物质证据，不能由分类推定实际生活。账本提示的研究线索为宫廷政治、制度运作与中央—地方信息链研究。事件所处条件为同治无嗣使宗法与宫廷权力问题交织。后续变化应限于慈禧重新垂帘，晚清中枢进入新阶段。来源限制：以同治十三年（1874年）的同期文书为定位起点；官修记录、奏报或外交文本服务特定机构立场，具体日期、规模、地方执行与对手视角须以独立材料复核。
